\documentclass[11pt]{article}
\usepackage[margin=1in]{geometry}
\usepackage{amsmath,amssymb,amsthm}
\usepackage{graphicx}
\usepackage{hyperref}

\newtheorem{theorem}{Theorem}

\title{A Trace Observation on the Rank-Two Projection Sparsification Question}
\author{Automatic Math Discovery Smoke Test}
\date{June 12, 2026}

\begin{document}
\maketitle

\section*{Source and Extracted Question}

The source is Grigory Ivanov, \emph{No-dimensional results of combinatorial
convexity. Dimension strikes back}, arXiv:2602.20035. In Section 6, after a
deterministic Rudelson-type conjecture for positive semidefinite matrices, the
paper asks whether a linear number of summands can suffice under additional
structure; the displayed example asks whether \(k=Cd\) is possible when each
\(A_i\) is an orthogonal projection of rank \(2\).

\begin{center}
\includegraphics[width=\textwidth]{figures/open_problem_crop.png}
\end{center}

\section*{Claim}

Under the normalization stated immediately before the rank-two projection
question, namely
\[
  A_i\in \mathcal M_d,\qquad A_i\geq 0,\qquad
  \lambda_i\geq 0,\qquad \sum_{i=1}^n \lambda_i=1,\qquad
  \sum_{i=1}^n \lambda_i A_i = I_d,
\]
the case in which every \(A_i\) is an orthogonal projection of rank \(2\) is
vacuous in dimensions \(d\neq 2\). In dimension \(d=2\), the problem is solved
exactly with one summand.

\begin{theorem}
Let \(A_1,\ldots,A_n\in\mathcal M_d\) be orthogonal projections of rank \(2\).
If there are nonnegative scalars \(\lambda_1,\ldots,\lambda_n\) with
\(\sum_i\lambda_i=1\) and
\[
  \sum_{i=1}^n \lambda_i A_i = I_d,
\]
then \(d=2\). Consequently, for the rank-two projection formulation as printed,
there are no feasible instances in dimensions \(d>2\), and every feasible
instance in dimension \(2\) is represented exactly by a single summand.
\end{theorem}

\begin{proof}
The trace of an orthogonal projection is its rank. Hence
\(\operatorname{tr}(A_i)=2\) for every \(i\). Taking traces in the assumed
identity gives
\[
  d
  = \operatorname{tr}(I_d)
  = \operatorname{tr}\left(\sum_{i=1}^n \lambda_i A_i\right)
  = \sum_{i=1}^n \lambda_i \operatorname{tr}(A_i)
  = 2\sum_{i=1}^n \lambda_i
  = 2.
\]
Thus \(d=2\) is forced.

If \(d=2\), an orthogonal projection of rank \(2\) on the whole space is
\(I_2\). Pick any index \(i_1\) and set \(\beta_1=1\). Then
\[
  \left\|\beta_1 A_{i_1}-I_2\right\|_{S_\infty}=0.
\]
Thus \(k=1\) suffices in the only feasible dimension.
\end{proof}

\section*{Interpretation and Limitation}

This should be viewed as a normalization issue in the printed rank-two example,
not as a solution to a renormalized sparsification problem. A different
question, for example allowing the original convex coefficients to have total
\(d/2\), or replacing projections by positive semidefinite matrices normalized
to have trace \(d\), is not addressed by this packet.

\section*{Verification Notes}

No computation is needed. The entire argument is the trace identity above.
The main point for human review is to confirm that ``orthogonal projection of
rank \(2\)'' is meant in the standard sense, so that its trace equals \(2\).

\end{document}
